
\section{Clergy-Body or Congregational Freedom}

Editorial article in “Folkeladet” for May 18, 1898.

Many a time it seems almost hopeless to write on this subject, though for our Church it is so highly important. There are none deafer than the one who will not hear, and it seems futile to argue this matter with those pastors—of all Norwegian church bodies—who at any cost will have a clergy-body, even if they are willing to call it a congregational body or even a free church. It is not merely lack of sincerity that makes them so altogether unwilling to discuss the question of church order in a becoming manner; nor is it only the consciousness that something would have to be attempted for the sake of “the Body of Christ,” even if, in order to make room for a free and living congregation, they should have to relinquish an illegitimate power and authority, equally ruinous to both parties, according to the Word of God; “not as those who would rule over the congregations” (1 Pet. 5:3); for there is also an excuse for the state-church clerical party: from childhood onward, and by a state-church upbringing, they have been so altogether permeated with prejudices over against the congregational question, that even the simplest statements grounded in God’s Word concerning it sound like to them a wholly foreign and incomprehensible language, a language for which many of them lack the spiritual prerequisites to understand, even if they were willing.

It is with them as it was with the priesthood in the time of the Reformers; they were so thoroughly worked through with notions of “the extraordinary holiness of the office and its unlimited mediating authority,” so terrifyingly secure because their views had the sanction of many, many centuries and the testimony of the fathers on their side, that the Reformers’ word concerning the individual’s justification solely by his own personal faith sounded to them like a foreign and heretical speech, a speech sprung from self-interest, lust for domination, and ignobility, however much it was grounded—indeed precisely because, passing by the fathers, it was grounded in God’s Word itself.

It ought surely to be a simple and indisputable truth that the ordering of the congregation and the mutual connection of the congregations, upon which their spiritual life and activity so greatly depend, must be grounded and regulated by God’s Word just as much as doctrine. And nevertheless it seems almost impossible to get the state-church clerical party of The United Church to understand or acknowledge this plain truth, and on that basis to ground the discussion apart from all personal and foreign considerations.

And if it were not because men and women round about in the congregations in all the bodies had begun to think seriously and independently about this most holy thing among us, the Body of Christ, the congregation, it would be a thankless labor to convince narrow-sighted, prejudice-bound pastors of former days.

In “Lutheraneren” for May 11 there stand, on the occasion of Prof. Briggs, who, although judged by the Presbytery to be a heretic, is nevertheless to continue as theological teacher at Union Seminary, the following tendentious remarks:

“This shows how helpless a church body can become over against false doctrine, when it does not itself own and administer its seminary. It may declare a theological professor to be a heretic, and yet he continues all the same to spread his heresy within the body. It is precisely this relation that the leading Augsburgers have wished to place Augsburg Seminary in toward The United Church, and now most recently toward the Free Church.”

For the present we shall call attention to the fact that in the remark cited above—consistently with the language usage of The United Church—the expression “church body” is used, not “the congregations.” “Church body” is a word or concept that is not found in the New Testament; but in its historical development it has always denoted a communal connection in which the pastors have been either sole rulers or nearly sole rulers, partly on the basis of a pretended divine authority, as in the Catholic Church and the state churches, partly on the basis of a gradually heightened supremacy with the prescriptive right of age, as in most church bodies outside the state churches. The congregations and their right have in all these cases—except at the very most as wards—been left out of consideration.

Let us now make some observations drawn from historical reality concerning these “church bodies” and their “helplessness” over against heretical teachers.

When the Catholic church body, or “the Church,” as they were pleased to call it, deposed and condemned, and in part burned, such teachers as Wycliffe, Hus, and Jerome of Prague as heretics and false teachers, then surely on account of its power it was anything but “helpless.” But were the men named above therefore, from a Protestant point of view, heretics, or was there in the unrestricted power of the priesthood any guarantee for the congregations or for pure doctrine? Was not Luther entirely in his right when, on the ground of God’s Word and on behalf of God’s congregation, he protested against those teachers of theology whom the Catholic Church, with all authority and unanimity, appointed in all schools, and in whose appointment the people, who paid, had neither access nor right.

But, you say, they were Catholics and Jesuits, and we all know that they were heretical and jesuitical.

Let us then pass over to a church body with which we are well acquainted, the Norwegian Lutheran State Church. There too the priesthood had everything to say in matters of doctrine or theology, and the congregations were till the very latest unable to exercise any influence upon the choice of their own pastors or teachers.

In Hauge’s time the entire priesthood in Norway, together with the civil government, was rationalistic, and the teachers at the schools of higher learning were of the same sort. The state-church body, that is to say, the priesthood, had all the power; they were anything but “helpless,” and no one could prevent them from appointing rationalistic teachers of pastors. Is that the condition which the defenders of clergy-bodies regard as the one most in accord with God’s Word?

But, people say, that was a state church; here in this country we have free-church bodies. Very well, let us consider one of them, for example the Norwegian Synod. It makes claim to be free-church enough and asserts that it represents the congregations, in that they have the right to send representatives to the annual meetings. As is well known, the Norwegian Synod owns and controls its pastoral seminary as far as that is at all possible; otherwise the Anti-Missourians would doubtless have taken it too. Thus it is not helpless in this question. Are the congregations thereby now secured with respect to the doctrinal soundness of the theological teachers? [—] will be able to answer that better than Pastor Kildahl in Chicago.

For it is an equally well-known fact that, however much one may speak of free-church constitution and the right of the congregations in the Norwegian Synod, it is still the pastors who, partly through their pastors’ conferences, partly through the annual meetings, have the decisive, not to say the absolute, voice in all questions of the body, especially concerning doctrine and theological teachers.

If now the priesthood in the Norwegian Synod has the decisive voice in the choice of theological teachers, will it not require of them that they have the doctrinal position of the priesthood, or even of a majority of the priesthood? If now this position should happen to be false in doctrine, what remedy then have the congregations? Who then become helpless, the body or the congregations?

Is not this precisely what has happened in the Norwegian Synod, at least according to the Anti-Missourians’ understanding and course of action.

It helped nothing that the body owned and controlled the chair; it helped nothing that a number of congregations and pastors protested against the false doctrine of grace-election, and over it carried on the bitterest separations in earnest—the false teachers remained standing, upheld by the majority of the body, that is to say the majority of the pastors.

And what was the only remedy that the opposing part of the body, namely the Anti-Missourians, could find? Division, division.

The anti-Missourian pastors split the body; indeed, more than that, they went around from congregation to congregation, split as many as they could, brought bitterness and party-spirit into hearts, tore up the Norwegian church people anew under the leadership of a German professor, and led them into litigious wrangling and desire for property not rightfully theirs, which has followed the Anti-Missourians down to their latest Augsburg lawsuit.

The facts lie before everyone’s eyes, and each may draw his own conclusions. In truth, the clergy-bodies have been neither a spiritual nor a temporal blessing for our Norwegian church people. When the priesthood with its leaders has gained so complete an upper hand in The United Church, where by the constitution one had nevertheless hoped to secure to the congregations some right, that they have even been able to cast out of the body congregations and lawfully chosen delegates partly in order to satisfy personal revenge, partly because the congregations had maintained a school which The United Church now in court claims that the expelled members of the congregation have always lawfully been corporate members of—what security then have the congregations in such a clergy-body, either for the doctrine of a seminary or for its property?

When that same church body has in addition incorporated itself, and in many ways gathered considerable property, so that the priesthood can gradually feel itself fairly secure in its power, what regard will one then have for the congregations except as a taxable subject?

. . . Oh no, take freedom and voluntariness away from the congregations, and put a priesthood with clerical rule in their stead—however well it may begin, it will soon end in an absolute and spiritless clerical dominion. And of The United Church it cannot even be said that it has the advantage of beginning well.

The Presbyterians have a remedy against Prof. Briggs’s rationalism; they can withdraw their support from his school.

And a free congregation has this remedy against a pastoral seminary which does not uphold the congregation’s right and does not labor for the preservation of spiritual life: it can withhold its support. And if it is free and unhindered by all bonds of church-body organization and chains of the clergy, it can, without any regard to such things, lay its powers there where God’s Word is taught purely and rightly, and the life of God is esteemed and nurtured. . . .

S. Oftedal
